\begin{problem}[2019年第36届全国中学生物理竞赛复赛][旅行车圆凳茶杯稳定问题]{104}
如图a，旅行车上有一个半径为 $R$ 的三脚圆凳(可视为刚性结构)，三个相同凳脚的端点连线（均水平）构成边长为 $a$ 的等边三角形，凳子质心位于其轴上的 $G$ 点．半径为 $r$ 的一圆筒形薄壁茶杯放在凳面上，杯底中心位于凳面中心 $O$ 点处，茶杯质量为 $m$（远小于凳子质量），其中杯底质量为 $\frac{m}{5}$ （杯壁和杯底各自的质量分布都是均匀的），杯高为 $H$（与杯高相比，杯底厚度可忽略）．杯中盛有茶水，茶水密度为 $\rho$ ．重力加速度大小为 $g$．
\begin{subproblem}
为了使茶水杯所盛茶水尽可能多并保持足够稳定，杯中茶水的最佳高度是多少？
\end{subproblem}
\begin{subproblem}
现该茶水杯（杯中茶水高度最佳）的底面边缘刚好缓慢滑移到与圆凳的边缘内切于 $D$ 点时静止（凳面边有小凸缘，可防止物体滑出；凳面和凳面边的凸缘各自的质量分布都是均匀的），且 $OD\perp AC$（见图b），求此时旅行车内底板对各凳脚的支持力相对于滑移前（该茶水杯位于凳面中心处）的改变.
\end{subproblem}
\end{problem}